% Research fragment.  Requires amsmath, amssymb, amsthm and hyperref.
% Theorem/proposition/corollary/remark environments are supplied by the parent.
% All claims proved here; global novelty is not claimed.
\section{A Gamma coefficient calculus for shifted harmonic sums}
\label{harm:sec:coefficient-calculus}

The exact families in this section complement the manuscript's Hurwitz-jet
calculus.  They involve a nonlinear but structured class of generalized
harmonic-number numerators and admit a finite reduction for every rational
denominator that decays quadratically.  The unshifted specialization is
classical: it recovers the Aomoto--Drinfeld--Zagier generating function for
height-one multiple zeta values.  Young's fractional-order Barnes--Hurwitz
generating function \cite{harm:Young2024}, and the height-one regularization
formulas recalled by Kaneko and Sakata \cite{harm:KanekoSakata}, are directly
relevant prior work.  Our purpose is a proved shifted-parameter calculus,
with a precise two-variable explanation of its finite parts.  We do not
claim that the kernel, its unshifted specialization, or its classical
low-depth Euler sums are new.

Throughout this section, $\Re a>0$ and $\log(n+a)$ is the principal
logarithm.  Define the elementary harmonic coefficients by
\begin{equation}\label{harm:eq:elementary}
 \prod_{k=1}^{n}\left(1+\frac{y}{k}\right)
 =\sum_{r=0}^{n} E_{n,r}y^r,
 \qquad E_{0,0}=1,
 \qquad E_{n,r}=0\quad(r>n).
\end{equation}
Thus, with $H_n^{(j)}=\sum_{k=1}^n k^{-j}$ and $H_n=H_n^{(1)}$,
\begin{align}
 E_{n,0}&=1, &E_{n,1}&=H_n,\\
 E_{n,2}&=\frac{H_n^2-H_n^{(2)}}2,
 &E_{n,3}&=\frac{H_n^3-3H_nH_n^{(2)}+2H_n^{(3)}}6.
 \label{harm:eq:first-elementaries}
\end{align}
Equivalently, $E_{n,r}$ is the sum of $1/(k_1\cdots k_r)$ over
$1\le k_1<\cdots<k_r\le n$.  In particular,
$E_{n,r}=O_r((1+\log(n+1))^r)$.

The Gamma coefficients to which these sums reduce are
\begin{equation}\label{harm:eq:gamma-coefficients}
 F(a,y):=\frac{\Gamma(a)\Gamma(1-y)}{\Gamma(a-y)}
       =\sum_{k=0}^{\infty}C_k(a)y^k.
\end{equation}
This is analytic for sufficiently small $y$, and $C_0(a)=1$.  Put
\begin{equation}\label{harm:eq:generalized-harmonics}
 A_1(a)=\gamma+\psi(a),\qquad
 A_j(a)=\zeta(j)-\zeta(j,a)\quad(j\ge2).
\end{equation}
The Gamma Taylor series gives
\begin{equation}\label{harm:eq:bell-generator}
 F(a,y)=\exp\left(\sum_{j=1}^{\infty}\frac{A_j(a)}j y^j\right),
 \qquad
 kC_k(a)=\sum_{j=1}^{k} A_j(a)C_{k-j}(a).
\end{equation}
Consequently every $C_k$ is a polynomial with rational coefficients in
generalized harmonic functions.  For a positive integer $a=N$,
$A_j(N)=H_{N-1}^{(j)}$, and $C_k(N)$ is the complete homogeneous
harmonic sum of depth $k$ on $1,\ldots,N-1$.

For compact formulas, write
\begin{equation}\label{harm:eq:ABCD}
 A=\gamma+\psi(a),\quad B=\zeta(2)-\psi'(a),\quad
 C=\zeta(3)+\frac{\psi''(a)}2,\quad
 D=\zeta(4)-\frac{\psi^{(3)}(a)}6.
\end{equation}
Then
\begin{align}
 C_1(a)&=A,\\
 C_2(a)&=\frac{A^2+B}{2},\\
 C_3(a)&=\frac{A^3+3AB+2C}{6},\\
 C_4(a)&=\frac{A^4+6A^2B+3B^2+8AC+6D}{24}.
 \label{harm:eq:first-C}
\end{align}

\subsection{An exact generating theorem}

Define
\begin{equation}\label{harm:eq:shifted-series}
 Z_r(s,a):=\sum_{n=0}^{\infty}\frac{E_{n,r}}{(n+a)^s},
 \qquad \Re s>1.
\end{equation}
For $a=1$, the convention with the largest index first gives
$Z_r(s,1)=\zeta(s,\{1\}^r)$.

\begin{theorem}[Shifted harmonic Gamma generating function]
\label{harm:thm:gamma-generating}
For $|x|<\Re a$ and $|y|<1$,
\begin{equation}\label{harm:eq:double-generator}
 \sum_{m=1}^{\infty}\sum_{r=0}^{\infty}
 Z_r(m+1,a)x^m y^{r+1}
 =F(a,y)-F(a-x,y).
\end{equation}
In particular, for all integers $r\ge0$ and $m\ge1$,
\begin{equation}\label{harm:eq:coefficient-reduction}
 \boxed{\displaystyle
 Z_r(m+1,a)=\frac{(-1)^{m+1}}{m!}
                 \frac{d^m}{da^m}C_{r+1}(a).}
\end{equation}
For $\Re a,\Re b>0$, the endpoint difference is the convergent identity
\begin{equation}\label{harm:eq:resolvent-difference}
 \boxed{\displaystyle
 \sum_{n=0}^{\infty}E_{n,r}
 \left(\frac1{n+a}-\frac1{n+b}\right)
 =C_{r+1}(b)-C_{r+1}(a).}
\end{equation}
\end{theorem}

\begin{proof}
The ordinary generating function of the elementary harmonic coefficients
follows from the binomial theorem:
\begin{equation}\label{harm:eq:ordinary-generator}
 \sum_{n=0}^{\infty} E_{n,r}u^n
 =\frac{(-\log(1-u))^r}{r!(1-u)},\qquad |u|<1.
\end{equation}
Indeed,
$\sum_n(1+y)_n u^n/n!=(1-u)^{-1-y}$, and extraction of $y^r$
gives the assertion.  The integrable difference in
\eqref{harm:eq:resolvent-difference} therefore has generating function
\begin{align}
 &\sum_{r\ge0}y^r\sum_{n\ge0}E_{n,r}
       \left(\frac1{n+a}-\frac1{n+b}\right)\notag\\
 &\qquad=\int_0^1
       \bigl(u^{a-1}-u^{b-1}\bigr)(1-u)^{-1-y}\,du
       =B(a,-y)-B(b,-y)
       =\frac{F(b,y)-F(a,y)}y.
 \label{harm:eq:beta-difference}
\end{align}
The difference integral converges for $\Re y<1$: its numerator vanishes
to first order at $u=1$.  Equality with the difference of beta functions
first holds for $\Re y<0$ and extends by holomorphy, with the singularity
at $y=0$ removable.  Near $u=0$, $\Re a,\Re b>0$ provides an integrable
majorant.  For the sums, absolute and locally uniform convergence follow
from $\sum_r E_{n,r}|y|^r=(1+|y|)_n/n!=O(n^{|y|})$ and the quadratic
decay of the resolvent difference.  Coefficient extraction proves
\eqref{harm:eq:resolvent-difference}.

Apply the same identity with $b=a$ and the first parameter $a-x$.
Since
\[
 \frac1{n+a-x}-\frac1{n+a}
 =\sum_{m=1}^{\infty}\frac{x^m}{(n+a)^{m+1}},
\]
the claimed double generating function follows.  The imposed bound on
$x$ justifies this expansion uniformly in $n$, and coefficient extraction
from $F(a-x,y)$ gives \eqref{harm:eq:coefficient-reduction}.
\end{proof}

At $a=1$, $F(1,y)=1$, so \eqref{harm:eq:double-generator} becomes
\[
 \sum_{m\ge1,r\ge0}\zeta(m+1,\{1\}^r)x^m y^{r+1}
 =1-\frac{\Gamma(1-x)\Gamma(1-y)}{\Gamma(1-x-y)},
\]
the classical height-one generating function.  In this specialization
the right side is symmetric in $x,y$, giving the standard height-one
duality.  There is generally no such symmetry for an arbitrary shift $a$.

\subsection{Explicit polygamma reductions and antiderivatives}

Set $P_j=\psi^{(j)}(a)$, and retain $C_0=1$ and the explicit $C_k$ in
\eqref{harm:eq:first-C}.  The first two denominator powers, at depths
zero through three, are as follows:
\begin{align}
 Z_0(2,a)&=P_1,\\
 Z_1(2,a)&=C_1P_1-\frac{P_2}{2},\\
 Z_2(2,a)&=C_2P_1-\frac{C_1P_2}{2}+\frac{P_3}{6},\\
 Z_3(2,a)&=C_3P_1-\frac{C_2P_2}{2}
                    +\frac{C_1P_3}{6}-\frac{P_4}{24};
 \label{harm:eq:square-table}
\end{align}
and
\begin{align}
 Z_0(3,a)&=-\frac{P_2}{2},\\
 Z_1(3,a)&=-\frac12\left(C_1P_2+P_1^2-\frac{P_3}{2}\right),\\
 Z_2(3,a)&=-\frac12\left[
 C_2P_2+C_1\left(P_1^2-\frac{P_3}{2}\right)
       +\frac{P_4}{6}-P_1P_2\right],\\
 Z_3(3,a)&=-\frac12\left[
 C_3P_2+C_2\left(P_1^2-\frac{P_3}{2}\right)
 +C_1\left(\frac{P_4}{6}-P_1P_2\right)
 -\frac{P_5}{24}+\frac{P_1P_3}{3}+\frac{P_2^2}{4}\right].
 \label{harm:eq:cube-table}
\end{align}
These are finite closed forms, not recursively defined Euler sums.
For example,
\begin{equation}\label{harm:eq:linear-example}
 \sum_{n=0}^{\infty}\frac{H_n}{(n+a)^2}
 =\bigl(\psi(a)+\gamma\bigr)\psi'(a)-\frac{\psi''(a)}2.
\end{equation}
Each entry is an explicit antiderivative formula as well:
\begin{equation}\label{harm:eq:antiderivative}
 \int Z_r(2,a)\,da=C_{r+1}(a)+\mathrm{constant},\qquad
 \int Z_r(m+1,a)\,da=-\frac1mZ_r(m,a)+\mathrm{constant}
 \quad(m\ge2).
\end{equation}
The first formula identifies a finite polynomial in polygammas as the
primitive; it avoids leaving an unevaluated harmonic Dirichlet series.

For a useful nonintegral specialization, put $L=\log2$.  The half-point
values $\zeta(j,1/2)=(2^j-1)\zeta(j)$ give
\begin{align}
 Z_1(2,1/2)&=7\zeta(3)-6L\zeta(2),\\
 Z_2(2,1/2)&=\frac{15}{2}\zeta(4)-14L\zeta(3)+6L^2\zeta(2),\\
 Z_3(2,1/2)&=31\zeta(5)-13\zeta(2)\zeta(3)-15L\zeta(4)
                       +14L^2\zeta(3)-4L^3\zeta(2).
 \label{harm:eq:half-examples}
\end{align}
Thus the last identity, written without $E_{n,3}$, is
\begin{align}
 &\sum_{n=0}^{\infty}
 \frac{H_n^3-3H_nH_n^{(2)}+2H_n^{(3)}}{(2n+1)^2}\notag\\
 &\quad=\frac32\left[
 31\zeta(5)-13\zeta(2)\zeta(3)-15L\zeta(4)
 +14L^2\zeta(3)-4L^3\zeta(2)\right].
 \label{harm:eq:half-cubic-unpacked}
\end{align}
The factor $3/2$ is $3!/4$, from the harmonic normalization and the
conversion $(n+1/2)^{-2}=4(2n+1)^{-2}$.

\subsection{Closure under arbitrary rational kernels}

\begin{corollary}[Rational-kernel reduction]\label{harm:cor:rational}
Let $R(z)$ be rational, have all its finite poles among $-a_1,\ldots,-a_J$
with $\Re a_j>0$, and satisfy $R(z)=O(z^{-2})$ at infinity.  Write
\[
 R(z)=\sum_{j=1}^{J}\sum_{\ell=1}^{L_j}
              \frac{c_{j,\ell}}{(z+a_j)^\ell}.
\]
Then $\sum_jc_{j,1}=0$, and
\begin{align}\label{harm:eq:rational-compiler}
 \sum_{n=0}^{\infty}E_{n,r}R(n)
 ={}&-\sum_{j=1}^{J}c_{j,1}C_{r+1}(a_j)\notag\\
 &+\sum_{j=1}^{J}\sum_{\ell=2}^{L_j}
 \frac{(-1)^\ell c_{j,\ell}}{(\ell-1)!}
           C_{r+1}^{(\ell-1)}(a_j).
\end{align}
\end{corollary}
\begin{proof}
The vanishing sum of residues is the coefficient of $z^{-1}$ at infinity.
It permits replacing each simple-pole summand by its difference from a
fixed reference resolvent.  Apply \eqref{harm:eq:resolvent-difference} to
these differences and \eqref{harm:eq:coefficient-reduction} to all repeated
poles.  All combined series are absolutely convergent.
\end{proof}

For real $a>0$ and real $b\ne0$, a compact example is
\begin{equation}\label{harm:eq:quadratic-kernel}
 \boxed{\displaystyle
 \sum_{n=0}^{\infty}\frac{E_{n,r}}{(n+a)^2+b^2}
 =\frac{C_{r+1}(a+ib)-C_{r+1}(a-ib)}{2ib}
 =\frac{\Im C_{r+1}(a+ib)}b.}
\end{equation}
The limit $b\to0$ recovers $Z_r(2,a)=C_{r+1}'(a)$.

\subsection{The moving pole and all-depth finite parts}

The noninteger order of the Barnes--Hurwitz generating kernel is the
right variable for explaining why the endpoint finite parts are also
Gamma coefficients.  This is more delicate than substituting $s=1$ into
the divergent defining series.

\begin{theorem}[Two-variable pole removal]\label{harm:thm:moving-pole}
For $\Re a>0$, initially define
\begin{equation}\label{harm:eq:Barnes-kernel}
 \mathcal B(s,y;a)=\sum_{n=0}^{\infty}
            \frac{(1+y)_n}{n!(n+a)^s},
 \qquad \Re s>1+\Re y.
\end{equation}
Near $(s,y)=(1,0)$ there is a jointly holomorphic function
$\mathcal R(s,y;a)$ such that
\begin{equation}\label{harm:eq:moving-pole}
 \Gamma(s)\mathcal B(s,y;a)
 =\frac1{s-1-y}+\mathcal R(s,y;a),
 \qquad
 \mathcal R(1,y;a)=\frac{1-F(a,y)}y.
\end{equation}
If
\begin{equation}\label{harm:eq:reciprocal-gamma}
 \frac1{\Gamma(1+z)}=\sum_{j=0}^{\infty}g_jz^j,
\end{equation}
then
\begin{equation}\label{harm:eq:all-depth-finite-part}
 \boxed{\displaystyle
 Z_r(1+z,a)=\sum_{j=0}^{r}\frac{g_j}{z^{r+1-j}}
          +g_{r+1}-C_{r+1}(a)+O(z).}
\end{equation}
The $O(z)$ term is locally uniform in $a$ on the right half-plane.
\end{theorem}

\begin{proof}
The Mellin representation, initially in a common half-plane where
absolute convergence justifies summation under the integral, is
\[
 \Gamma(s)\mathcal B(s,y;a)
 =\int_0^\infty t^{s-1}e^{-at}(1-e^{-t})^{-1-y}\,dt.
\]
Split it at $t=1$ and remove its leading small-$t$ term.  The remainder is
\begin{align}\label{harm:eq:holomorphic-remainder}
 \mathcal R(s,y;a)
 ={}&\int_0^1 t^{s-2-y}
       \left[e^{-at}\left(\frac{t}{1-e^{-t}}\right)^{1+y}-1\right]dt\notag\\
 &+\int_1^\infty t^{s-1}e^{-at}(1-e^{-t})^{-1-y}\,dt.
\end{align}
The bracket is $O(t)$ uniformly on compact subsets of $(y,a)$ with
$\Re a>0$.  Hence the first integral is holomorphic whenever
$\Re(s-y)>0$, and the second is holomorphic in $s,y$ because of
exponential decay.  Derivatives introduce only bounded powers of
$\log t$ near zero and polynomial factors at infinity; the same
majorants justify all mixed derivatives on compact subsets.  The
subtracted integral is $(s-1-y)^{-1}$.

For $\Re y<0$, substitute $u=e^{-t}$ at $s=1$ to obtain
\[
 \mathcal B(1,y;a)=B(a,-y)=-\frac{F(a,y)}y.
\]
Consequently $\mathcal R(1,y;a)=(1-F(a,y))/y$, and this equality extends
holomorphically across $y=0$.

Extracting $y^r$ first in a region $|y|<|s-1|$ gives
\[
 \Gamma(s)Z_r(s,a)=\frac1{(s-1)^{r+1}}+R_r(s,a),
 \quad R_r(s,a):=[y^r]\mathcal R(s,y;a),
 \quad R_r(1,a)=-C_{r+1}(a).
\]
Both sides extend meromorphically to a neighborhood of $s=1$.  Multiply
by the reciprocal Gamma Taylor series.  Its coefficient $g_{r+1}$
from the pole contributes to the constant term, in addition to
$R_r(1,a)$.  This proves \eqref{harm:eq:all-depth-finite-part}.
\end{proof}

The first reciprocal Gamma coefficients are
\begin{align}
 g_0&=1,\\
 g_1&=\gamma,\\
 g_2&=\frac{\gamma^2-\zeta(2)}2,\\
 g_3&=\frac{\gamma^3-3\gamma\zeta(2)+2\zeta(3)}6,\\
 g_4&=\frac{\gamma^4-6\gamma^2\zeta(2)+3\zeta(2)^2
                    +8\gamma\zeta(3)-6\zeta(4)}{24}.
 \label{harm:eq:first-g}
\end{align}
Thus the complete principal parts and constants for $r=0,1,2,3$ are
\begin{align}
 Z_0(1+z,a)&=z^{-1}+\gamma-A+O(z),\\
 Z_1(1+z,a)&=z^{-2}+\gamma z^{-1}
                   +g_2-\frac{A^2+B}{2}+O(z),\\
 Z_2(1+z,a)&=z^{-3}+\gamma z^{-2}+g_2z^{-1}
                   +g_3-\frac{A^3+3AB+2C}{6}+O(z),\\
 Z_3(1+z,a)&=z^{-4}+\gamma z^{-3}+g_2z^{-2}+g_3z^{-1}\notag\\
 &\hspace{25mm}+g_4-
       \frac{A^4+6A^2B+3B^2+8AC+6D}{24}+O(z).
 \label{harm:eq:finite-part-table}
\end{align}
For $a=1$, $C_k(1)=0$ for every $k\ge1$, and the constant is simply
$g_{r+1}$, consistently with height-one multiple-zeta regularization.

\subsection{Stieltjes constants, the literature comparison, and limits of the reduction}

Use the conventional normalization
\begin{equation}\label{harm:eq:Stieltjes-convention}
 \zeta(1+z,a)=\frac1z+
        \sum_{m=0}^{\infty}\frac{(-1)^m\gamma_m(a)}{m!}z^m.
\end{equation}
The $r=0$ finite-part formula says
$\gamma_0(a)=\gamma-C_1(a)=-\psi(a)$, as required.  More generally let
\[
 \Gamma(1+z)=\sum_{j\ge0}q_jz^j,
 \qquad R_{r,k}(a)=[y^rz^k]\mathcal R(1+z,y;a).
\]
The ordinary Stieltjes layer is exactly
\begin{equation}\label{harm:eq:Stieltjes-boundary}
 R_{0,k}(a)=q_{k+1}+
       \sum_{j=0}^{k}q_{k-j}\frac{(-1)^j\gamma_j(a)}{j!}.
\end{equation}
For every $r\ge0$ the coefficient of $z^m$, $m\ge0$, in the full
Laurent expansion of $Z_r(1+z,a)$ is
\begin{equation}\label{harm:eq:higher-coefficient}
 d_{r,m}(a)=g_{r+m+1}+
                  \sum_{k=0}^{m}g_{m-k}R_{r,k}(a).
\end{equation}
Our finite-part theorem determines $R_{r,0}(a)=-C_{r+1}(a)$.  It does
\emph{not} determine $R_{r,k}$ for $r,k\ge1$ from Gamma coefficients.
Those mixed Taylor coefficients are higher harmonic Stieltjes data.
Any claim that all of them reduce to polygamma values would require
additional identities.

The distinction between two shifted conventions matters.  Karg\i n,
Dil, Cenkci, and Can \cite{harm:Kargin2023} study
\[
 \zeta_H(s,a)=\sum_{n=0}^{\infty}
       \frac{\sum_{k=0}^{n}(k+a)^{-1}}{(n+a)^s}.
\]
For general $a$, its harmonic numerator is shifted along with its
denominator, whereas ours is the unshifted $H_n$.  The two should not be
identified.  At $a=1$ they satisfy
\begin{equation}\label{harm:eq:harmonic-literature-comparison}
 Z_1(s,1)=\zeta_H(s,1)-\zeta(s+1).
\end{equation}
The known constant term
$\gamma_H(0,1)=(\gamma^2+\zeta(2))/2$, quoted in
\cite{harm:Kargin2023}, therefore gives
$\operatorname{FP}_{s=1}Z_1(s,1)=(\gamma^2-\zeta(2))/2=g_2$,
exactly as above.  Their higher harmonic Stieltjes constants, and
Young's higher-depth analogues \cite{harm:Young2023,harm:Young2024},
are the appropriate comparison for the still nontrivial mixed
coefficients.  Coffey's nonlinear harmonic reductions
\cite{harm:Coffey2011} already include familiar specializations such as
$Z_1(2,1)=\zeta(3)$ and $Z_2(2,1)=\zeta(4)$.

An entirely convergent integral representation of the mixed data is
available directly from \eqref{harm:eq:holomorphic-remainder}.  Put
\[
 K_r(t,a)=\frac{e^{-at}[-\log(1-e^{-t})]^r}
                  {r!(1-e^{-t})}.
\]
Then
\begin{align}\label{harm:eq:mixed-integrals}
 R_{r,k}(a)=\frac1{k!}\bigg\{
 &\int_0^1(\log t)^k
   \left[K_r(t,a)-\frac{(-\log t)^r}{r!t}\right]dt\notag\\
 &+\int_1^\infty(\log t)^kK_r(t,a)\,dt\bigg\}.
\end{align}
For example, $R_{0,1}(a)$ combines with $R_{0,0}(a)$ by
\eqref{harm:eq:higher-coefficient} to give $-\gamma_1(a)$.  Formula
\eqref{harm:eq:mixed-integrals} is a useful exact definition and a
numerical verification route, but is not asserted to be a closed
evaluation of all mixed data.

